\item Dos espejos rectangulares planos, ambos perpendiculares a una hoja de papel horizontal, se colocan borde a borde, con sus superficies reflectoras perpendiculares entre sí.
\begin{enumerate}
    \item Un rayo de luz en el plano del papel incide en uno de los espejos a un ángulo arbitrario de incidencia $\theta_1$. Demuestre que la dirección final del rayo, después de la reflexión desde ambos espejos, es opuesta a su dirección inicial.
    \item \textbf{¿Qué pasaría si?} El papel es sustituido por un tercer espejo plano que toca los bordes de los otros dos y se coloca perpendicular a ambos, creando un retroreflector de esquina de cubo. Un rayo de luz incide desde cualquier dirección dentro del octante del espacio limitado por las superficies reflectoras. Demuestre que el rayo se reflejará una vez desde cada espejo y que su dirección final será opuesta a su dirección original.
\end{enumerate}